% Protokół badania — etap 1 (około jednej strony)
% Kompilacja: pdflatex protokol.tex (TeX Live; dokument po polsku, babel + T1)
\documentclass[10pt,a4paper]{article}

\usepackage[utf8]{inputenc}
\usepackage[T1]{fontenc}
\usepackage[polish]{babel}
\usepackage[margin=1.5cm]{geometry}
\usepackage{booktabs}
\usepackage{amsmath}
\usepackage{enumitem}
\usepackage{hyperref}
\hypersetup{colorlinks=true, linkcolor=black, urlcolor=blue, citecolor=black}
\setlist{nosep, leftmargin=1.4em}
\pagestyle{empty}
\setlength{\parskip}{0.2em}
\setlength{\parindent}{0pt}

\newcommand{\repoURL}{https://github.com/FylyppS/rozpoznawanie-wzorcow}
% zwarty nagłówek sekcji (bez odstępów \section*)
\newcommand{\hd}[1]{\par\vspace{0.55em}{\normalsize\bfseries #1}\par\vspace{0.1em}}

\begin{document}
\small

{\large\bfseries Protokół badania --- etap 1: prezentacja tematu}\\
{\small Rozpoznawanie wzorców, studia II stopnia, 2026/2027, Grupa 2 \hfill Dariusz Salzburg (103931), Filip Sejud (103928) \hfill 29.10.2026}

\vspace{0.3em}
\textbf{Temat 06. Obciążenie i zmienność oceny przez walidację krzyżową.} Publikacja: R.~Kohavi,
\emph{A Study of Cross-Validation and Bootstrap for Accuracy Estimation and Model Selection}, IJCAI 1995.
Repozytorium (kod, protokół, surowe wyniki): \url{\repoURL}

\hd{1. Pytania badawcze i hipotezy}
\begin{itemize}
  \item \textbf{P1.} Jak liczba części $k$ walidacji krzyżowej wpływa na obciążenie i zmienność szacowanego błędu?
    \emph{H1:} $\hat e_{CV}$ jest obciążone pesymistycznie (zawyża błąd), najsilniej dla $k=2$ (trening na $n/2$ obiektów),
    a przy $k=10$ obciążenie jest bliskie zera; zmienność jest największa dla $k=2$.
    \emph{H1b:} efekt jest silniejszy dla niestabilnego drzewa CART niż dla stabilnego GaussianNB.
  \item \textbf{P2.} Co zmienia stratyfikacja przy klasie rzadkiej ($\approx 10\%$)? \emph{H2:} stratyfikacja zmniejsza
    zmienność $\hat e_{CV}$ i eliminuje foldy bez obiektów klasy rzadkiej; wpływ na obciążenie jest niewielki.
\end{itemize}

\hd{2. Wielkości mierzone}
Jednostką powtórzenia jest niezależny zbiór treningowy $D_r$ ($r=1,\dots,R$), nie pojedynczy fold.
Dla modelu dopasowanego na $D_r$ i $k$-krotnej CV o foldach $D_{r,1},\dots,D_{r,k}$:
\[
\hat e_{CV,r} = \frac{1}{k}\sum_{j=1}^{k} e_j, \qquad
\text{bias}_k \approx \frac{1}{R}\sum_{r=1}^{R}\bigl(\hat e_{CV,r} - e_{test,r}\bigr), \qquad
\text{zmienność}_k = \operatorname{sd}_r\bigl(\hat e_{CV,r} - e_{test,r}\bigr),
\]
gdzie $e_j$ to odsetek pomyłek modelu uczonego na $D_r\setminus D_{r,j}$ i ocenianego na $D_{r,j}$,
a $e_{test,r}$ --- odsetek pomyłek modelu dopasowanego na \emph{całym} $D_r$ i ocenionego na dużym, niezależnym
zbiorze testowym $T$ (ocena referencyjna, przybliżenie błędu prawdziwego). Przy równych foldach średnia błędów foldów
równa się ilorazowi liczby pomyłek i $n$ (wzór Kohaviego); zapisujemy obie wartości. Podajemy też $\operatorname{sd}_r \hat e_{CV,r}$,
błąd standardowy obciążenia $\operatorname{sd}/\sqrt{R}$, a w eksperymencie własnym czułość klasy rzadkiej i liczbę foldów bez tej klasy.

\hd{3. Generator danych (utrwalony)}
$X\mid y \sim \mathcal N(\mu_y, I_{10})$, $\mu_1=-\mu_0=\tfrac{\delta}{2}u$, $u=(1,1,1,1,1,0,0,0,0,0)/\sqrt5$
(5 cech informacyjnych, 5 szumowych), $\delta=2$ (odległość Mahalanobisa między średnimi klas), $y\sim\text{Bernoulli}(\pi_1)$.
Błąd Bayesa: $\Phi(-\delta/2)=15{,}9\%$ dla $\pi_1=0{,}5$; dla $\pi_1=0{,}1$ próg $t=\ln 9/\delta$, błąd Bayesa $7{,}0\%$,
czułość klasy rzadkiej $46\%$. Rozkład jest zgodny z założeniami GaussianNB (model dobrze wyspecyfikowany, stabilny),
a granica ukośna względem osi jest trudna dla CART (model niestabilny).

\hd{4. Odtworzenie (adaptacja)}
$R=30$ zbiorów $D_r$ po $n=200$ obiektów (ziarno $1000+r$); jeden zbiór $T$, $N=10\,000$ (ziarno $90001$, niezależne od $D_r$).
Modele: \texttt{GaussianNB()} i CART \texttt{DecisionTreeClassifier(random\_state=0)} (bez przycinania, parametry domyślne).
$k\in\{2,5,10\}$, \texttt{KFold(shuffle=True, random\_state=}$2000+r$\texttt{)} --- te same podziały dla obu modeli.
\begin{center}\small
\begin{tabular}{@{}lll@{}}
\toprule
 & Kohavi (1995) & Nasza adaptacja \\
\midrule
Dane & 7 zbiorów UCI (+ losowy) & generator syntetyczny o znanym błędzie Bayesa \\
Modele & C4.5, Naive-Bayes (MLC++) & CART, GaussianNB (scikit-learn 1.8.0) \\
Punkt odniesienia & holdout na reszcie zbioru, 50 powtórzeń & duży niezależny test $N=10\,000$, 30 powtórzeń \\
Metody & $k\in\{2,5,10,20\}$, leave-$k$-out, bootstrap .632 & $k\in\{2,5,10\}$; bez bootstrapu \\
Miara & trafność & odsetek pomyłek (= 1 $-$ trafność) \\
Stratyfikacja & wariant przy klasach zbalansowanych & osobny eksperyment z klasą rzadką 10\% \\
\bottomrule
\end{tabular}
\end{center}

\hd{5. Eksperyment własny}
$\pi_1=0{,}1$ (ok.\ 20 obiektów klasy rzadkiej na 200), $R=30$, ziarna $3000+r$ / $4000+r$, test $T'$ ($N=10\,000$, ziarno $90002$).
Na \emph{tym samym} $D_r$: \texttt{KFold(5)} i \texttt{StratifiedKFold(5)} z tym samym ziarnem. Porównujemy obciążenie, zmienność,
czułość klasy rzadkiej oraz liczbę foldów bez klasy rzadkiej (czułość nieokreślona --- zapisywana jako brak wyniku z powodem).

\hd{6. Kontrola poprawności i budżet}
\begin{itemize}
  \item Duży test służy wyłącznie do oceny estymatora; nie dobieramy według niego parametrów modeli ani generatora.
  \item Foldy nie są niezależne: zbiory uczące sąsiednich foldów dzielą $(k-2)/(k-1)$ obiektów; niepewność liczymy po $r$, nie po foldach.
  \item Ziarna, parametry, metryki i podziały są utrwalone w \texttt{cvbias/protocol.py}; każde uruchomienie zapisuje wyniki per fold
    (CSV), skład klas w foldach i \texttt{meta.json} (wersje bibliotek, liczba dopasowań, czas). Zmiany protokołu odnotowujemy w raporcie.
  \item Budżet: odtworzenie $30\cdot2\cdot(2+5+10+1)=1080$ dopasowań, eksperyment własny $30\cdot2\cdot(5+5+1)=660$,
    rozdzielenie wpływu rozmiaru treningu od podziału (etap 3; model na $n(k-1)/k$ obiektach) $180$ --- razem $1920\le2500$.
    Próba uruchomienia (3 zbiory): 174 dopasowania w $0{,}16$ s (CPU).
\end{itemize}

\end{document}
